\section{Vector-Valued Quantization and a Threshold at Eight Bits}
\label{sec:multibit-extension}

Theorem~\ref{thm:CK} motivates the vector-valued question of whether a
\(k\)-bit summary is optimized by retaining \(k\) input coordinates. The
answer is negative in general. This section gives explicit counterexamples,
describes their coding-theoretic mechanism, and isolates the range that
remains open.

The outcome is unexpectedly sharp. Shortened Hamming quantizers violate
the coordinate benchmark for \(k=8,9,10\), and hence, by adding independent
coordinate outputs, for every \(k\geq 8\). The same construction misses the
benchmark at \(k=7\) by only \(1/64\) in the high-noise coefficient. A broad
collection of searches below eight bits also failed to produce a violation.
This leads us to conjecture that the coordinate projection is optimal for
\(2\leq k\leq7\).

Throughout this section, entropies and mutual informations are measured in
bits. Let
\begin{align*}
 X&\sim\operatorname{Unif}(\mathbb F_2^n),\qquad Y=X\oplus Z,\\
 Z_i&\stackrel{\mathrm{iid}}{\sim}\operatorname{Bernoulli}(p),
 \qquad 0\leq p\leq\frac12.
\end{align*}
and write
\[
 h_2(p)=-p\log_2p-(1-p)\log_2(1-p),
 \qquad C(p)=1-h_2(p).
\]

\subsection{Problem Formulation and Known Results}

For \(n\geq k\), define
\begin{equation}
 M_k(n,p)
 :=\max_{f:\mathbb F_2^n\to\mathbb F_2^k} I(f(X);Y).
 \label{eq:multibit-Mkn}
\end{equation}
The \(k\)-coordinate projection
\[
 D_k(x)=(x_1,\ldots,x_k)
\]
satisfies \(I(D_k(X);Y)=kC(p)\). We call \(D_k\), and any version obtained
by permuting input coordinates and relabeling its \(2^k\) outputs, a
\emph{dictator tuple}. The most direct vector extension of the
Courtade--Kumar assertion is therefore
\begin{equation}
 M_k(n,p)\stackrel{?}{\leq} kC(p)
 \qquad\text{for every }n\geq k.
 \label{eq:vector-CK}
\end{equation}

There are three useful landmarks in the literature.

\begin{enumerate}
\item For \(k=1\), inequality~\eqref{eq:vector-CK} is the original
Courtade--Kumar conjecture~\cite{CK}, proved in this paper.

\item Huleihel and Ordentlich proved the complementary endpoint: for every
\(f:\mathbb F_2^n\to\mathbb F_2^{n-1}\),
\[
 I(f(X);Y)\leq(n-1)C(p).
\]
Their theorem also extends from the BSC to every binary-input memoryless
output-symmetric channel~\cite{HO}. Notice that this is an assertion with
\(k=n-1\), not an assertion for a fixed \(k\) and arbitrary \(n\); hence it
does not conflict with the fixed-\(k\) counterexamples below.

\item Chandar and Tchamkerten formulated the general optimization
problem~\eqref{eq:multibit-Mkn} and characterized its positive-rate
asymptotics~\cite{CT}. If \(k/n\to R\in(0,1)\), then
\begin{equation}
 \begin{split}
 \frac{M_k(n,p)}{n}&\longrightarrow
 1-h_2\!\left(p\star h_2^{-1}(1-R)\right),\\
 p\star q&=p(1-q)+(1-p)q.
 \end{split}
 \label{eq:CT-asymptotic}
\end{equation}
This limit can exceed \(RC(p)\), so the vector assertion cannot hold at
every positive rate. At finite blocklength, they exhibited violations
from the perfect \([15,11,3]\) Hamming code and the \([23,12,7]\) Golay
code. They also observed that the \([7,4,3]\) Hamming code does not
violate the bound and explicitly asked for constructions with \(k<11\).
\end{enumerate}

One elementary point prevents a possible misunderstanding: a linear
\emph{code} can produce a nonlinear \emph{quantizer}.

\begin{proposition}[Affine maps obey the coordinate bound]
\label{prop:affine-multibit}
If \(f(x)=Ax+b\) is affine, with output rank \(r\leq k\), then
\[
 I(f(X);Y)\leq rC(p)\leq kC(p).
\]
\end{proposition}

\begin{proof}
It is enough to treat a rank-\(r\) linear map. Since \(X\) is uniform,
\(AX\) is uniform on \(\mathbb F_2^r\). Moreover, \(Z=X\oplus Y\) is
independent of \(Y\), and hence
\[
 I(AX;Y)=r-H(AZ).
\]
Choose \(r\) columns of \(A\) that form an invertible submatrix and call
their index set \(J\). Conditioning on \(Z_{J^c}\), the variable \(AZ\)
is a translate of an invertible image of \(Z_J\). Consequently,
\[
 H(AZ)\geq H(AZ\mid Z_{J^c})=H(Z_J)=r h_2(p),
\]
which proves the claim.
\end{proof}

Thus every counterexample must use genuine nonlinearity. The constructions
below use linear codebooks, but the nearest-codeword map into those
codebooks is nonlinear.

\subsection{Coset Quantizers and the High-Noise Expansion}
\label{subsec:coset-diagnostic}

Let \(C\leq\mathbb F_2^n\) be a binary linear code of dimension \(k\),
with full-rank parity-check matrix \(H\in\mathbb F_2^{(n-k)\times n}\).
For each syndrome \(a\in\mathbb F_2^{n-k}\), choose a minimum-weight
coset leader \(e_a\) satisfying \(He_a=a\), with \(e_0=0\), and define
\begin{equation}
 Q(x)=x\oplus e_{Hx}\in C.
 \label{eq:coset-quantizer}
\end{equation}
After composing \(Q\) with any bijection from \(C\) to
\(\mathbb F_2^k\), this is a \(k\)-bit quantizer. The choice among tied
minimum-weight leaders is part of the construction.

If \(X\) is uniform, the change of variables
\[
 X=Q(X)\oplus E,
 \qquad E=e_{HX},
\]
makes \(Q(X)\) uniform on \(C\), makes \(E\) uniform on the set of coset
leaders, and makes the two variables independent. Since \(Y\) is uniform
on \(\mathbb F_2^n\), translation invariance gives the exact identity
\begin{equation}
 I(Q(X);Y)=n-H(E\oplus Z).
 \label{eq:coset-information}
\end{equation}
This reduces the entire information calculation to the entropy of a noisy
coset leader.

Set \(\rho=1-2p\) and define
\begin{equation}
 B(E):=\sum_{j=1}^n
 \left(\mathbb E[(-1)^{E_j}]\right)^2.
 \label{eq:leader-B}
\end{equation}
For a set \(S\), the Walsh coefficient of \(E\oplus Z\) is
\[
 \rho^{|S|}\mathbb E\!\left[(-1)^{\sum_{j\in S}E_j}\right].
\]
Expanding relative entropy from the uniform distribution at \(\rho=0\)
therefore yields
\begin{align}
 I(Q(X);Y)
 &=\frac{\rho^2}{2\ln2}B(E)+O(\rho^4),
 \label{eq:code-high-noise}\\
 kC(p)
 &=\frac{\rho^2}{2\ln2}k+O(\rho^4).
 \label{eq:benchmark-high-noise}
\end{align}
In particular,
\begin{equation}
 \begin{split}
 B(E)>k\quad\Longrightarrow\quad I(Q(X);Y)>kC(p)\\
 \text{for all \(p<1/2\) sufficiently close to \(1/2\)}.
 \end{split}
 \label{eq:B-counterexample}
\end{equation}
This exact second-order test is the main design criterion used below.

\subsection{Shortened-Hamming Counterexamples for
\texorpdfstring{\(k=8,9,10\)}{k=8,9,10}}
\label{subsec:shortened-hamming}

Identify the nonzero vectors of \(\mathbb F_2^4\) with the integers
\(1,\ldots,15\) through their binary expansions. For
\(s\in\{1,2,3\}\), let \(H_s\) have as its columns
\[
 \{s+1,s+2,\ldots,15\}.
\]
Thus \(C_s=\ker H_s\) is the shortened Hamming code
\begin{equation}
 C_s=[15-s,11-s,3],
 \qquad n=15-s,
 \qquad k=11-s.
 \label{eq:shortened-Hamming-parameters}
\end{equation}

The retained syndrome \(8\) will be an anchor. Choose the following
minimum-weight leader for every syndrome \(a\in\mathbb F_2^4\):
\begin{equation}
 e_a=
 \begin{cases}
 0, & a=0,\\
 \mathbf e_a, & a\in\{s+1,\ldots,15\},\\
 \mathbf e_8\oplus\mathbf e_{8\oplus a}, & a\in\{1,\ldots,s\}.
 \end{cases}
 \label{eq:star-leaders}
\end{equation}
Here \(\mathbf e_a\) denotes the unit vector in the coordinate whose
parity-check column is \(a\). Both \(8\) and \(8\oplus a\) are retained.
The last line has syndrome \(a\), and it is minimum-weight because the
column \(a\) itself was removed. The missing syndromes therefore use
weight-two leaders that form a star around the common anchor \(8\).

The leader set consists of one word of weight zero, \(n\) words of weight
one, and \(s\) words of weight two. More importantly, the anchor occurs
in \(s+1\) leaders, each of the \(s\) leaves occurs in two leaders, and
each of the remaining \(14-2s\) coordinates occurs in one leader. Hence
\begin{align}
 B_s
 &=\left(1-\frac{s+1}{8}\right)^2
   +s\left(1-\frac{2}{8}\right)^2\notag\\
 &\quad +(14-2s)\left(1-\frac{1}{8}\right)^2\notag\\
 &=\frac{s^2-76s+735}{64},
 \label{eq:star-B}\\
 B_s-k
 &=\frac{s^2-12s+31}{64}.
 \label{eq:star-B-minus-k}
\end{align}

\begin{proposition}[Counterexamples below eleven bits]
\label{prop:k8-k10-counterexamples}
For each \(k\in\{8,9,10\}\), there are an \(n=k+4\), a quantizer
\(f:\mathbb F_2^n\to\mathbb F_2^k\), and an open interval of crossover
probabilities ending at \(1/2\) on which
\[
 I(f(X);Y)>kC(p).
\]
\end{proposition}

\begin{proof}
Use the quantizer~\eqref{eq:coset-quantizer} with the code and leaders in
\eqref{eq:shortened-Hamming-parameters}--\eqref{eq:star-leaders}. For
\(s=1,2,3\), corresponding respectively to \(k=10,9,8\),
\[
 B_s-k\in\left\{\frac5{16},\frac{11}{64},\frac1{16}\right\}>0.
\]
The conclusion follows from the high-noise expansion
\eqref{eq:B-counterexample}.
\end{proof}

\begin{proposition}[Failure for every \(k\geq8\)]
\label{cor:failure-all-k-ge-8}
For every integer \(k\geq8\), inequality~\eqref{eq:vector-CK} fails for
some input dimension and some \(p\in(0,1/2)\).
\end{proposition}

\begin{proof}
Start with the \(8\)-bit counterexample \(f_8\). For \(t=k-8\), append
\(t\) fresh independent input coordinates and define
\[
 F(x,u)=(f_8(x),u),
 \qquad u\in\mathbb F_2^t.
\]
The two channel blocks are independent, so
\[
 I(F(X,U);Y_X,Y_U)=I(f_8(X);Y_X)+tC(p)>kC(p).
\]
\end{proof}

The proof is analytic and uses only the sign of the exact rational number
\(B_s-k\). Direct entropy evaluation gives a more quantitative picture.
For \(y\in\mathbb F_2^n\), we computed
\begin{equation}
 \Pr(E\oplus Z=y)
 =\frac1{16}\sum_{a\in\mathbb F_2^4}
 p^{d_H(y,e_a)}(1-p)^{n-d_H(y,e_a)}
 \label{eq:direct-leader-convolution}
\end{equation}
and substituted the resulting distribution in
\eqref{eq:coset-information}. In Table~\ref{tab:shortened-hamming}, the
leader enumerator is
\(\sum_a z^{\operatorname{wt}(e_a)}\). The positive intervals and maxima
are numerical; the code parameters, leader enumerators, and \(B_s-k\)
are exact.

\begin{table*}[t]
\caption{Shortened-Hamming counterexamples. The code parameters,
leader enumerators, and high-noise excesses are exact; the intervals and
maxima are obtained by direct entropy evaluation.}
\label{tab:shortened-hamming}
\centering
{\footnotesize\setlength{\tabcolsep}{5pt}\renewcommand{\arraystretch}{1.12}
\begin{tabular}{@{}cccccc@{}}
\toprule
\(k\) & code & leader enumerator & \(B_s-k\) & numerical \(p\)-interval
& \(\max_p\Delta_k(p)/k\) \\ \midrule
10 & \([14,10,3]\) & \(1+14z+z^2\) & \(5/16\)
& \((0.088658,1/2)\) & \(4.48440\!\times\!10^{-3}\) \\
9 & \([13,9,3]\) & \(1+13z+2z^2\) & \(11/64\)
& \((0.152276,1/2)\) & \(1.83320\!\times\!10^{-3}\) \\
8 & \([12,8,3]\) & \(1+12z+3z^2\) & \(1/16\)
& \((0.265563,1/2)\) & \(3.21308\!\times\!10^{-4}\) \\
\bottomrule
\end{tabular}}
\end{table*}
Here
\[
 \Delta_k(p)=I(f(X);Y)-kC(p).
\]
The maximizing crossover probabilities are approximately \(0.199318\),
\(0.248809\), and \(0.332782\), respectively. Simple interior witness
points are
\begin{align*}
 \Delta_{10}(0.2)/10&=0.0044842878,\\
 \Delta_9(0.25)/9&=0.0018330159,\\
 \Delta_8(1/3)/8&=0.0003212940.
\end{align*}
Direct convolution and an independent Walsh-transform calculation agreed
to within \(10^{-14}\) at the tested crossover probabilities.

\begin{remark}[Order of shortening and decoding]
The successful operation is to remove columns from the Hamming
parity-check matrix and then choose a new minimum-weight syndrome table.
Deleting output bits after decoding with the full \([15,11,3]\) Hamming
quantizer is a different operation and did not reproduce these gains.
The shared anchor in~\eqref{eq:star-leaders} is also essential: it aligns
the coordinate biases that enter \(B(E)\).
\end{remark}

\subsection{Searches Below Eight Bits}
\label{subsec:below-eight-searches}

The star construction itself explains why eight is a natural boundary.
For \(s=4,5,6\), corresponding to \(k=7,6,5\), formula
\eqref{eq:star-B-minus-k} gives
\begin{table}[t]
\caption{High-noise coefficients of the shortened-Hamming family below
eight output bits.}
\label{tab:hamming-below-eight}
\centering
{\scriptsize\setlength{\tabcolsep}{3.5pt}\renewcommand{\arraystretch}{1.08}
\begin{tabular}{@{}ccccc@{}}
\toprule
\(k\) & \(n\) & leader enumerator & \(B_s\) & \(B_s-k\) \\ \midrule
7 & 11 & \(1+11z+4z^2\) & \(447/64\) & \(-1/64\) \\
6 & 10 & \(1+10z+5z^2\) & \(95/16\) & \(-1/16\) \\
5 & 9  & \(1+9z+6z^2\)  & \(315/64\) & \(-5/64\) \\
\bottomrule
\end{tabular}}
\end{table}
The seven-bit case is strikingly close: its leading high-noise term misses
the dictator tuple by only \(1/64\). A dense numerical evaluation found no
positive gap on \(0<p<1/2\); for example, the per-bit deficits at
\(p=0.1,0.2,0.4\) are respectively
\[
 0.0115176,\qquad 0.00410469,\qquad 0.000105113.
\]

We then tested several ways in which the Hamming near-miss might be
improved. Table~\ref{tab:k7-searches} records the scope of each computation and
the best high-noise statistic found. A negative-curve outcome means that
the full mutual-information curve was also evaluated on a dense grid and
at its numerical stationary candidates; it is not an analytic
impossibility proof.

\begin{table*}[t]
\caption{Computational searches for seven-bit code quantizers. A
``negative curve'' records dense numerical evaluation and stationary-point
checks, not an impossibility theorem for the corresponding family.}
\label{tab:k7-searches}
\centering
{\footnotesize\setlength{\tabcolsep}{4pt}\renewcommand{\arraystretch}{1.13}
\begin{tabular}{@{}
  >{\raggedright\arraybackslash}p{0.28\linewidth}
  >{\raggedright\arraybackslash}p{0.36\linewidth}
  >{\centering\arraybackslash}p{0.13\linewidth}
  >{\raggedright\arraybackslash}p{0.15\linewidth}@{}}
\toprule
Family & Search performed & best \(B\) & outcome \\ \midrule
Shortened Hamming, \(k=7\)
& All \(\binom{15}{4}=1365\) shortenings and all \(451{,}920\)
minimum-distance tie tables
& \(447/64\) & negative curve \\

Primitive BCH \([15,7,5]\)
& All \(1941\) puncturings deleting \(0,\ldots,4\) coordinates, with
multistart optimization of minimum-distance ties
& \(447/64\) & best case returns the Hamming near-miss \\

First- and second-order Reed--Muller family
& All \(651\) seven-dimensional codes
\(\mathrm{RM}(1,4)\subset C\subset\mathrm{RM}(2,4)\)
& \(6.850952\) & negative curves \\

Shortened binary Golay
& \(203\) sampled shortenings to dimension seven, with tie optimization
& \(6.895\) & negative curves \\

Random parity-check matrices
& \(345{,}000\) distinct-column searches with redundancies \(5,6,7,8\)
& \(447/64\) & no improvement over Hamming \\

Linear quotients of the \(k=8\) counterexample
& All \(255\) rank-seven quotients of the \([12,8,3]\) quantizer
& \(6.28125\) & negative curves \\
\bottomrule
\end{tabular}}
\end{table*}

We also tested repeated parity-check columns, nonminimum coset
representatives, the extended \([16,7,6]\) BCH construction, the
\([15,5,7]\) BCH code, and \(\mathrm{RM}(1,4)=[16,5,8]\). Repeated-column
searches occasionally reached \(B=7\), but only through a degeneracy whose
entire information curve equals that of seven independent dictators; no
strict gain appeared. Multistart searches over arbitrary coset
representatives likewise reached equality but not a violation.

These computations support a structural claim about the code-based
mechanism, not a theorem about all quantizers. In particular, the exact
exhaustion of Hamming tie tables certifies the best \(B\) only within that
family. A quantizer with a smaller high-noise coefficient could still,
in principle, cross the benchmark at an intermediate \(p\).

\subsection{Random Coding at \texorpdfstring{\(n=15,\ k=7\)}{n=15, k=7}}
\label{subsec:random-coding-k7}

The failed structured searches raise a natural question: is an unstructured
codebook more likely to find the missing gain? We tested the finite random
ensemble suggested by the rate-distortion picture. Choose a codebook
\(\mathcal C\) uniformly without replacement from all \(M=2^7\)-subsets of
\(\mathbb F_2^{15}\), map each input to a nearest codeword, and break each
tie independently and uniformly, once and for all.

Some geometric statistics of this ensemble are exact. Let
\[
 N=2^{15},\qquad
 S_d=\binom{15}{d},\qquad
 V_d=\sum_{i=0}^d S_i.
\]
For a fixed source word, let \(D\) be its distance to the random codebook
and \(T\) the number of codewords at that minimum distance. Then
\begin{align}
 \Pr(D>t)
 &=\frac{\binom{N-V_t}{M}}{\binom NM},
 \label{eq:random-distance-tail}\\
 \Pr(D=d,T=m)
 &=\frac{\binom{S_d}{m}\binom{N-V_d}{M-m}}{\binom NM}.
 \label{eq:random-distance-ties}
\end{align}
For \(n=15\) and \(M=128\), these formulas give
\begin{align*}
 \mathbb E D&=2.6608532,& \frac{\mathbb E D}{15}&=0.1773902,\\
 \Pr(T>1)&=0.4945889,& \mathbb E T&=2.1208972.
\end{align*}
Thus almost half of the source words have a nearest-neighbor tie.

For an arbitrary quantizer \(U=Q(X)\), not necessarily a coset quantizer,
the high-noise coefficient is
\begin{equation}
 B(Q)=\sum_{j=1}^n\sum_u\Pr(U=u)
 \left(\mathbb E[(-1)^{X_j}\mid U=u]\right)^2,
 \label{eq:general-B}
\end{equation}
and
\begin{equation}
 I(U;Y)=\frac{\rho^2}{2\ln2}B(Q)+O(\rho^4).
 \label{eq:general-high-noise}
\end{equation}
We enumerated the full cube for \(128\) independently sampled codebooks
and evaluated every reported mutual information exactly for that sampled
quantizer. Thus the only Monte Carlo error is across codebooks, not within
a codebook. The comparison below includes a structured code of the same
length and size, and the shortened-Hamming near-miss of a different length.
For the BCH code, the displayed syndrome table was optimized over its
minimum-distance ties.

\begin{table*}[t]
\caption{Random and structured seven-bit quantizers. Mutual information
is evaluated exactly for each sampled quantizer; only the averaging over
random codebooks introduces Monte Carlo uncertainty.}
\label{tab:random-versus-structured}
\centering
{\footnotesize\setlength{\tabcolsep}{8pt}\renewcommand{\arraystretch}{1.12}
\begin{tabular}{@{}lccc@{}}
\toprule
Quantity
& random \((15,2^7)\), mean
& BCH \([15,7,5]\)
& Hamming \([11,7,3]\) \\ \midrule
\(H(U)\)
& \(6.975169\)
& \(7\)
& \(7\) \\
\(B(Q)\)
& \(6.379431\)
& \(6.850952\)
& \(6.984375\) \\
\(\Delta_7(0.1)/7\)
& \(-0.0505444\)
& \(-0.0294541\)
& \(-0.0115176\) \\
\(\Delta_7(0.4)/7\)
& \(-0.00262873\)
& \(-0.000697266\)
& \(-0.000105113\) \\
\bottomrule
\end{tabular}}
\end{table*}

Across the \(128\) random codebooks, \(B(Q)\) ranged only from
\(6.342219\) to \(6.423068\), and no sampled codebook beat the coordinate
benchmark at any tested crossover probability. The random Voronoi cells
were also appreciably unbalanced: their mean cell-size standard deviation
was \(47.31\) around the ideal size \(256\). This explains the loss
\(H(U)<7\) already at \(p=0\). The sample standard deviation of \(B(Q)\)
was \(0.01856\), giving the normal-approximation \(95\%\) interval
\([6.3762,6.3826]\) for its ensemble mean.

Because ties are common, we separately optimized their assignments for
the high-noise statistic. On \(64\) additional codebooks, a multistart
coordinate-ascent heuristic increased the mean \(B\) from \(6.3809\) to
\(6.6954\), with a best value \(6.71365\). This is a material improvement,
but it remains well below both \(7\) and the shortened-Hamming value
\(6.984375\). Since the tie optimization is heuristic, it is evidence
rather than an ensemble upper bound.

The experiment suggests that short-block structure matters in this
problem. Translation-equivariant code quantizers have equal-size cells
and can align their coordinate biases; a generic random Voronoi partition
has neither feature. This conclusion is qualitatively consonant with the
recent finite-block covering results of Riasat and Mahdavifar, where
structured Reed--Solomon-based constructions often outperform random
codebooks in average Hamming covering radius~\cite{RMcover}. The two
statements should not be conflated: their objective is average
nearest-codeword distance, whereas ours is BSC mutual information. The
agreement is a shared finite-block phenomenon, not a formal implication.

\subsection{A Seven-Bit Threshold Conjecture}
\label{subsec:seven-bit-conjecture}

The validity of the coordinate bound is monotone in the number of retained
bits in a useful sense. If a \(k\)-bit counterexample exists, appending
one fresh dictator gives a \((k+1)\)-bit counterexample. Conversely, if
the bound is valid for \(K\) bits in every dimension, then it is valid for
every \(k<K\): append \(K-k\) fresh dictators to a hypothetical \(k\)-bit
counterexample. Thus the set of fixed output sizes for which the bound is
universally true must be an initial segment of the positive integers.

The constructions above show that this initial segment ends no later than
seven. We conjecture that it ends exactly there.

\begin{conjecture}[Seven-bit threshold]
\label{conj:seven-bit-threshold}
For every \(2\leq k\leq7\), every \(n\geq k\), every
\(f:\mathbb F_2^n\to\mathbb F_2^k\), and every \(0\leq p\leq1/2\),
\begin{equation}
 I(f(X);Y)\leq k\bigl(1-h_2(p)\bigr).
 \label{eq:seven-bit-conjecture}
\end{equation}
Equality is attained by a \(k\)-coordinate projection.
\end{conjecture}

By the padding observation, it would be enough to prove
Conjecture~\ref{conj:seven-bit-threshold} for \(k=7\). Several pieces of
evidence point in the same direction.

\begin{enumerate}
\item \emph{The Hamming family changes sign at exactly the proposed
threshold.} Its exact excess coefficient is
\((s^2-12s+31)/64\): it is \(1/16\) at \(k=8\) and \(-1/64\) at \(k=7\).
This is not a numerical accident or a small interior-\(p\) margin.

\item \emph{Linear maps cannot be counterexamples.} By
Proposition~\ref{prop:affine-multibit}, a violation must combine nonlinear
cells with enough symmetry to preserve almost all \(k\) output bits.
Our unsuccessful postprocessing experiments illustrate the tension:
coarsening an eight-bit counterexample destroys the bias alignment much
faster than it reduces the benchmark.

\item \emph{The first high-noise target is concrete.} A proof of
\begin{equation}
 B(Q)\leq7
 \qquad\text{for every }Q:\mathbb F_2^n\to\mathbb F_2^7
 \label{eq:seven-bit-B-target}
\end{equation}
would settle the sharp second-order comparison at \(p=1/2\). Every
seven-bit family we tested satisfies this inequality, with the shortened
Hamming construction coming closest among the nontrivial examples.

\item \emph{Random coding points away from a hidden generic
counterexample.} Random nearest-neighbor quantizers are substantially
worse than the structured near-misses, even after aggressive tie
optimization. If the conjecture is false, the counterexample is likely to
require a rather special multiway partition rather than merely a good
covering code.
\end{enumerate}

There is also a clear reason the scalar theorem does not immediately
iterate. Writing \(U=(U_1,\ldots,U_k)\), the chain rule gives
\[
 I(U;Y)=\sum_{i=1}^k I(U_i;Y\mid U_1,\ldots,U_{i-1}),
\]
but after conditioning on the previous output bits, \(X\) is uniform only
on a cell of the induced partition, not on the full cube. The scalar
Courtade--Kumar theorem therefore does not apply to the conditional terms.
A proof of the seven-bit conjecture will likely need either a conditional
version of the scalar inequality or a genuinely multiway entropy-flow
argument.

\begin{remark}[A possible route for \(k=7\)]
The conjecture for all \(2\leq k\leq7\) reduces to one endpoint, \(k=7\).
A plausible program is to prove the high-noise inequality
\eqref{eq:seven-bit-B-target}, identify its equality cases, obtain a
multiway low-noise boundary inequality for balanced \(128\)-cell
partitions, and then bridge the two channel regimes. The \(1/64\) Hamming
near-miss provides a useful stress test for every proposed estimate.
\end{remark}
